\chapter{Четырёхтактная резонансная петля и предел открытой устойчивости}

Перед итоговой заморозкой был проведён обратный аудит ранней онтологии
проекта. В исходных текстах материя описывалась как свет, задержанный в
четыре раза и замкнутый в скрученную тороидальную форму. Поток должен был
входить во внутренний канал, проходить его и выходить с другой стороны,
сохраняя частицу как самосогласованный резонанс.

Эта формулировка не была строгой теорией скорости. Однако поздний проект
действительно нашёл фазу порядка четыре
\begin{equation}
 1\longmapsto i\longmapsto-1\longmapsto-i\longmapsto1.
 \label{eq:v5-order-four-phase-history}
\end{equation}
После двух тактов возникает мёбиусный знак, а после четырёх состояние
возвращается полностью. Поэтому раннюю гипотезу можно проверить без
буквального утверждения $v=c/4$: локальная скорость остаётся световой, а
наблюдаемое замедление понимается как время пребывания во внутреннем
четырёхтактном маршруте.

\section{Археология исходной идеи}

В раннем тексте \path{архив-2025-2026/2025-12-истоки/habr/print1.md}
сформулирован образ ``материя --- это свет, замедленный в четыре раза и
закрученный в топологические формы''. В следующем тексте
\path{архив-2025-2026/2025-12-истоки/habr/print2.md} устойчивость связывается
с потоком, который ныряет во внутренний канал, выходит с обратной стороны
и замыкает цикл; мёбиусность означает возвращение состояния только после
повторного обхода.

Поздняя дедуктивная ветвь правильно понизила буквальный статус этих
образов. Тор, лента Мёбиуса, расслоение Хопфа и
$\mathbb{RP}^3\times S^1$ не являются одним и тем же пространством.
Строгим остатком мёбиусной картины стали двойное накрытие, спиновый знак и
нетривиальная $\mathbb Z_2$-голономия.

Новая подсказка состоит в том, что квадратный корень этой
$\mathbb Z_2$-линии имеет характер $\pm i$. Следовательно,
\eqref{eq:v5-order-four-phase-history} является точной математической
версией четырёхтактного возврата. Но прежний гейт источника корня также
доказал, что эта линия не назначена обязательным образом текущему пакету
$H_{15}$.

\section{Минимальный унитарный резонатор}

Рассмотрим внешний канал и внутреннюю петлю, соединённые одним
безпотерьным узлом. Пусть $r\in[0,1]$ есть амплитуда остаться во внутреннем
канале после обхода, а
\begin{equation}
 t=\sqrt{1-r^2}
\end{equation}
--- амплитуда обмена с внешним каналом. Полный внутренний возврат требует
четырёх локальных тактов, поэтому
\begin{equation}
 q(\omega)=e^{4i\omega}.
\end{equation}

После суммирования всех повторных обходов проходная амплитуда равна
\begin{equation}
 S(\omega)=\frac{r-q(\omega)}{1-rq(\omega)}.
 \label{eq:v5-ring-scattering}
\end{equation}
Для $|q|=1$ выполняется
\begin{equation}
 |S(\omega)|=1,
\end{equation}
то есть модель точно унитарна и не теряет энергию.

Фазовая задержка имеет вид
\begin{equation}
 \tau(\theta)=
 4\frac{1-r^2}{1+r^2-2r\cos\theta},
 \qquad \theta=4\omega.
 \label{eq:v5-ring-group-delay}
\end{equation}
В резонансе
\begin{equation}
 \tau_{\rm res}=4\frac{1+r}{1-r}.
 \label{eq:v5-ring-resonant-delay}
\end{equation}

\begin{proposition}[Точный четырёхтактный проход]
При $r=0$ и $t=1$ имеем $S=-q$. Входящая волна полностью направляется во
внутренний маршрут, проходит четыре локальных такта и полностью выходит
обратно. Эффективная задержка равна четырём, хотя локальная скорость на
каждом ребре не изменяется.
\end{proposition}

Это является строгой кинематической реализацией ранней фразы о
``замедленном в четыре раза свете''. Но при $r=0$ возбуждение не хранится:
оно покидает петлю после одного полного цикла.

\section{Трилемма устойчивости и открытости}

После выключения внешнего входа внутренняя амплитуда умножается на $r$ за
каждый полный обход. Вероятность сохранения после $n$ обходов равна
\begin{equation}
 P_n=r^{2n}.
 \label{eq:v5-ring-survival}
\end{equation}
Для любого $0\le r<1$ время жизни конечно. Бесконечная устойчивость
требует $r=1$, но тогда
\begin{equation}
 t=0,
\end{equation}
и петля полностью отсоединена от внешнего канала.

\begin{nogocriterion}[Запрет устойчивого открытого однопортового узла]
В минимальном взаимном однопортовом резонаторе нельзя одновременно иметь
ненулевую амплитуду входа-выхода и бесконечное время локального хранения.
Открытая петля является резонансом конечной ширины; строго связанная петля
не возбуждается через тот же порт.
\end{nogocriterion}

Этот результат является конечномерной версией общей границы между
резонансом и связанным состоянием в непрерывном спектре. Тёмное состояние
может быть защищено разрушительной интерференцией, но его амплитуда на
канале связи должна исчезнуть; см. \path{arXiv:2309.01501} и квантовые
графовые примеры \path{arXiv:1410.2846}.

\section{Проверка тёмных четырёхтактных мод}

Внутренний циклический сдвиг на четырёх состояниях имеет собственные
значения
\begin{equation}
 1,\quad i,\quad-1,\quad-i.
\end{equation}
Если внешний порт присоединён к одному внутреннему узлу, каждая из четырёх
фурье-мод имеет ненулевое перекрытие с портом. Тёмных состояний нет, и все
моды имеют конечную ширину.

Если порт связывается с равномерной суммой четырёх узлов, три моды
становятся точно тёмными. Однако их перекрытие с тем же внешним каналом
равно нулю. В силу взаимности они не могут ни свободно войти, ни свободно
выйти через этот порт. Слабое нарушение симметрии делает их доступными, но
одновременно превращает в долгоживущие, а не вечные резонансы.

\section{Граница движения частицы}

Узел связи можно поместить в любую точку однородного внешнего носителя.
Все такие положения изоспектральны. Это согласуется с возможностью
переносить дефект без изменения его внутреннего типа.

Но статическая вырожденность положений ещё не является законом движения.
Нужно вывести локальное правило, которое переносит сам узел связи с одного
ребра на следующее, сохраняя четырёхтактную фазу и накопленную амплитуду.
Текущий резонатор такого правила не содержит.

\section{Статус ранней гипотезы}

Обратный аудит даёт смешанный, но содержательный результат.

Положительно установлено:
\begin{enumerate}
 \item ``замедление'' можно понимать как задержку, не меняя локальный
 световой конус;
 \item фаза порядка четыре даёт точную последовательность
 $1,i,-1,-i,1$;
 \item унитарный внутренний маршрут может дать ровно четыре такта задержки;
 \item внешний наблюдатель видит единый рассеивающий объект, а не вещество,
 запертое в сосуде.
\end{enumerate}

Не установлено:
\begin{enumerate}
 \item обязательное назначение корневой линии пакету $H_{15}$;
 \item бесконечная устойчивость при свободном входе и выходе;
 \item динамический закон перемещения узла связи;
 \item связь резонансной задержки с наблюдаемой массой.
\end{enumerate}

\begin{theorem}[Резонансное прочтение без физического замыкания]
Ранняя идея света в скрученной петле допускает точную унитарную
реализацию: четырёхтактный внутренний маршрут создаёт эффективную задержку
без уменьшения локальной скорости распространения. Однако минимальная
однопортовая архитектура не совмещает свободный обмен с бесконечной
локализацией. Поэтому она является моделью резонансного прохождения, но не
завершённой моделью устойчивой движущейся частицы.
\end{theorem}

Следующая глава может подвести итог части II, сохранив резонансную петлю
как точную кинематическую подсказку и зафиксировав условия её возможного
переоткрытия: новый защищающий канал, невзаимный узел либо выведенная
динамика подвижного присоединения.

\begin{protocol}
\begin{enumerate}
 \item ранняя формулировка найдена: \textbf{да};
 \item точная фаза порядка четыре: \textbf{существует};
 \item локальная скорость $c/4$: \textbf{не утверждается};
 \item эффективная задержка в четыре такта: \textbf{получена};
 \item унитарный свободный вход-выход: \textbf{получен};
 \item бесконечная устойчивость при том же входе-выходе:
 \textbf{невозможна в минимальной модели};
 \item тёмная устойчивая мода: \textbf{условна и отсоединена};
 \item закон движения дефекта: \textbf{не выведен};
 \item физическое замыкание: \textbf{не достигнуто}.
\end{enumerate}
\end{protocol}

Машинный протокол сохранён в
\path{s2t_v5_order_four_resonant_loop_transport_gate_results.json}.